%% toc: P15 158 158.6,158.7
% Weak convergence by bounded continuous test functions on a topological space, point-mass boundary witnesses for the portmanteau inequalities (Durrett, Section 3.2, pp. 119--120), nice spaces and the extension theorem (Durrett, pp. 53--54), the cardinality of Borel sets (Durrett, p. 54), and the reading, error and density notation of Examples 3.2.5 and 3.2.6 (Durrett, pp. 117--118).
\chapter{Test Functions, Boundary Witnesses, Nice Spaces and Borel Cardinality}
\label{sec:p15-w02c-p15-004-weak-test-functions}

\Cref{sec:p15-weak-convergence-of-laws} defines weak convergence of
Borel probability laws on the line through distribution functions,
and \cref{thm:p15-weak-bounded-continuous} shows that the same relation
is obtained by testing against bounded continuous functions. Durrett
introduces that theorem as a characterization which, in his words,
``makes sense in any topological space''
\cite[p.~119]{DurrettPTEAuthor2019}. After the portmanteau theorem he
adds a remark built from point masses that approach the boundary of an
open set, to show which way the inequalities for open and closed sets
point \cite[p.~120]{DurrettPTEAuthor2019}. Earlier, he introduces
nice measurable spaces as the class on which Kolmogorov's extension
theorem generalizes, after stating that the theorem fails for arbitrary
measurable spaces \cite[p.~53]{DurrettPTEAuthor2019}; and after coding
Borel subsets of complete separable metric spaces by real numbers, he
records without proof that such sets are countable or of the cardinality
of the real line \cite[p.~54]{DurrettPTEAuthor2019}. This chapter makes
these statements exact and proves them, or proves exactly where they
fail. The source is Version~5 of
\emph{Probability: Theory and Examples}, dated 11 January 2019.

\emph{Purpose.} To state the test-function definition on an arbitrary
topological space, prove that every ingredient it uses exists there,
identify it with the distribution-function definition on the line, and
record what can and cannot be concluded from it without further
hypotheses on the space; to prove the source's boundary remark in
the generality in which it is stated; to relate Durrett's nice spaces
to the handbook's Borel real codings and test his nice-space form of the
extension theorem; to prove the cardinality statement for Borel
sets; and to record the source's reading of the birthday example, the
size of the error in its numerical approximation, and its notation for
densities in the central order statistic example.

\emph{Inputs.} A topological space and Borel probability measures on
it, or random elements with those laws, possibly defined on different
probability spaces; for the second part, an open set of real numbers,
a boundary point of it, and a real sequence in the set converging to
that point; for the third part, measurable spaces with one-to-one maps
into the line that are measurable in both directions, and consistent
finite-dimensional laws on their products; for the fourth part, a Borel
subset of a complete separable metric space; for the last part, the
birthday time and the scaled median of
\cref{thm:p15-weak-first-examples,thm:p15-weak-central-order-statistic}.

\emph{Conclusions.} The test-function relation is well defined on every
topological space; on the line it coincides with
\cref{def:p15-weak-law-convergence}; it is preserved by continuous maps;
its limits are unique on a metrizable space and need not be unique on a
space that is not. Point masses converge weakly exactly when their
locations converge, and every boundary approach to an open set gives
strict inequality in both portmanteau inequalities. A nice map is a
Borel real coding exactly when its image is Borel; consistent laws on
products of nice spaces have at most one extension, have exactly one when
the maps have Borel images, and, under the Axiom of Choice, need not have
any otherwise. A Borel subset of a
complete separable metric space is finite, countably infinite, or in
bijection with \(\mathbb R\). In the birthday example the limit
approximation falls short of the exact probability by between
\(1.71\) and \(1.74\) percent of it, and the source's symbol
\(P(Y_n=y)\) in the order statistic example names a density, not the
probability of an event.

Several statements go beyond the source text: the uniqueness of limits
on a metrizable space, the two-point counterexample to uniqueness, the
converse half of the point-mass characterization, the exact relation
between nice maps and Borel real codings, and the failure of the
extension for nice factors whose maps have non-Borel images. The source
asserts only that the test-function definition makes sense beyond the
line, states the boundary remark without proof, calls the nice-space
extension trivial without stating it, and gives no proof of the
cardinality statement.

\section{The test-function definition on a topological space}
\label{sec:p15-w02c-p15-004-topological}

\begin{hbdefinition}{def:p15-w02c-p15-004-topological-weak}{Bounded continuous functions and test-function weak convergence}
Let \((S,\mathcal T)\) be a topological space: \(\mathcal T\) is a
collection of subsets of \(S\), called open sets, that contains
\(\varnothing\) and \(S\) and is closed under finite intersections and
arbitrary unions. A set is closed when its complement is open. The
\emph{Borel sigma-algebra} is \(\mathcal B(S)\defeq\sigma(\mathcal T)\);
for \(S=\mathbb R\) with its usual topology this is
\(\mathcal B(\mathbb R)\) of \cref{def:p15-borel-sets}. A function
\(g:S\to\mathbb R\) is \emph{continuous} when \(g^{-1}(U)\in\mathcal T\)
for every open \(U\subseteq\mathbb R\), and
\[
 C_b(S)\defeq\{g:S\to\mathbb R:\ g\hbox{ continuous and }
               \sup_{x\in S}|g(x)|<\infty\}.
\]
For Borel probability measures \(\mu_n\) on \((S,\mathcal B(S))\),
\(n\in\{1,2,\ldots\}\cup\{\infty\}\), write
\(\mu_n\Rightarrow\mu_\infty\) when
\[
 \int_S g\,d\mu_n\longrightarrow\int_S g\,d\mu_\infty
 \qquad\hbox{for every }g\in C_b(S).
\]
For \(S\)-valued random elements \(X_n\) with laws \(\mu_n\)
(\cref{def:p15-pushforward-law}), \(X_n\Rightarrow X_\infty\) means
\(\mu_n\Rightarrow\mu_\infty\). For \(x\in S\), \(\delta_x\) is the
point mass at \(x\): \(\delta_x(B)=1\) if \(x\in B\) and \(0\)
otherwise, for \(B\in\mathcal B(S)\).
\index{weak convergence!on a topological space}
\index{bounded continuous functions}
\end{hbdefinition}

The definition uses the topology twice, to say which functions are
tested and which sets are measurable, and uses nothing else about \(S\).
In particular it needs no order, so the distribution-function definition
of \cref{def:p15-weak-law-convergence} has no counterpart in a general
space, while this one does.

\begin{hbtheorem}{thm:p15-w02c-p15-004-topological-weak}{proved=annex:p15-w02c-p15-004-topo}{The test-function definition is meaningful on every topological space}
Let \((S,\mathcal T)\) be a topological space.
\begin{enumerate}
\item Every continuous \(g:S\to\mathbb R\) is
\(\mathcal B(S)/\mathcal B(\mathbb R)\)-measurable. Every
\(g\in C_b(S)\) is integrable against every Borel probability measure
\(\mu\), with \(|\int g\,d\mu|\leq\sup_S|g|\). Hence each integral in
\cref{def:p15-w02c-p15-004-topological-weak} is a finite real number,
and for random elements on arbitrary probability spaces,
\(\mathbb E g(X_n)=\int g\,d\mu_n\).
\item For \(S=\mathbb R\) with the usual topology, the relation of
\cref{def:p15-w02c-p15-004-topological-weak} holds exactly when
\(\mu_n\Rightarrow\mu_\infty\) in the sense of
\cref{def:p15-weak-law-convergence}.
\item If \((T,\mathcal U)\) is a topological space,
\(h:S\to T\) is continuous (\(h^{-1}(V)\in\mathcal T\) for every
\(V\in\mathcal U\)), and \(\mu_n\Rightarrow\mu_\infty\) on \(S\), then
\(h\) is \(\mathcal B(S)/\mathcal B(T)\)-measurable and the image laws
\(\mu_n\circ h^{-1}\) satisfy
\(\mu_n\circ h^{-1}\Rightarrow\mu_\infty\circ h^{-1}\) on \(T\).
\item If the topology of \(S\) is induced by a metric \(d\), and
\(\mu_n\Rightarrow\mu\) and \(\mu_n\Rightarrow\nu\), then \(\mu=\nu\)
on \(\mathcal B(S)\).
\end{enumerate}
\end{hbtheorem}

Part~(iii) applied to a homeomorphism and its inverse shows that the
relation depends on \(S\) only up to homeomorphism. Part~(iv) is the
metric analogue of the uniqueness of a proper weak limit in
\cref{thm:p15-weak-quantile-foundations}. Its metric hypothesis cannot
simply be deleted.

\begin{hbtheorem}{thm:p15-w02c-p15-004-nonunique}{proved=annex:p15-w02c-p15-004-nonunique}{A two-point space on which weak limits are not unique}
Let \(S=\{a,b\}\) with \(a\neq b\) and
\(\mathcal T=\{\varnothing,\{a\},S\}\). Then \(\mathcal T\) is a
topology, \(\mathcal B(S)\) consists of all four subsets of \(S\), and
\(C_b(S)\) consists of the constant functions. Consequently every
sequence of Borel probability measures on \(S\) converges, in the sense
of \cref{def:p15-w02c-p15-004-topological-weak}, to every Borel
probability measure on \(S\). In particular the constant sequence
\(\delta_a\) converges both to \(\delta_a\) and to \(\delta_b\), although
\(\delta_a(\{a\})=1\neq0=\delta_b(\{a\})\). The topology \(\mathcal T\)
is not induced by any metric.
\end{hbtheorem}

\begin{hbwarning}{def:p15-w02c-p15-004-what-transfers}{What the general definition does not carry with it}
The source's claim is that the test-function condition is a meaningful
definition on any topological space; it is not a claim that the real-line
theory of \cref{sec:p15-weak-convergence-of-laws} holds there. The
quantile realization of \cref{thm:p15-weak-quantile-coupling}, the
distribution-function and ramp arguments, and the uniqueness of limits
all use structure of \(\mathbb R\). \Cref{thm:p15-w02c-p15-004-nonunique}
shows that uniqueness fails on a space whose continuous functions do not
separate points; \cref{thm:p15-w02c-p15-004-topological-weak}(iv)
restores it under a metric. No portmanteau or selection theorem on a
general space is asserted in this chapter.
\index{weak convergence!uniqueness of limits}
\end{hbwarning}

\section{Point masses and the boundary remark}
\label{sec:p15-w02c-p15-004-boundary}

\begin{hbtheorem}{thm:p15-w02c-p15-004-point-masses}{proved=annex:p15-w02c-p15-004-point-masses}{Point masses converge weakly exactly when their locations converge}
For real numbers \(x_n\) and \(x_\infty\),
\(\delta_{x_n}\Rightarrow\delta_{x_\infty}\) if and only if
\(x_n\to x_\infty\). Equivalently, if
\(\mathbb P(X_n=x_n)=1\) for \(n\in\{1,2,\ldots\}\cup\{\infty\}\), then
\(X_n\Rightarrow X_\infty\) exactly when \(x_n\to x_\infty\).
\end{hbtheorem}

\begin{hbtheorem}{thm:p15-w02c-p15-004-boundary-approach}{proved=annex:p15-w02c-p15-004-boundary-approach}{Durrett's remark after Theorem~3.2.11: approach to a boundary point}
Let \(G\subseteq\mathbb R\) be open and let \(x_\infty\in\partial G\),
with interior, closure and boundary as in
\cref{def:p15-weak-boundary-sets}.
\begin{enumerate}
\item \(x_\infty\notin G\), and there is a sequence \(x_n\in G\) with
\(x_n\to x_\infty\).
\item For any such sequence, let \(\mathbb P(X_n=x_n)=1\) for
\(n\in\{1,2,\ldots\}\cup\{\infty\}\). Then \(X_n\Rightarrow X_\infty\),
\(\mathbb P(X_n\in G)=1\) for every \(n\), and
\(\mathbb P(X_\infty\in G)=0\). Hence
\[
 \liminf_{n\to\infty}\mathbb P(X_n\in G)=1>0=\mathbb P(X_\infty\in G),
\]
so the inequality of \cref{thm:p15-weak-portmanteau}(ii) is strict
for \(G\).
\item The set \(K=G^c\) is closed, and
\[
 \limsup_{n\to\infty}\mathbb P(X_n\in K)=0<1=\mathbb P(X_\infty\in K),
\]
so the inequality of \cref{thm:p15-weak-portmanteau}(iii) is strict
for \(K\).
\item \(\partial K=\partial G\) and
\(\mathbb P(X_\infty\in\partial G)=1\). Neither \(G\) nor \(K\) is a
continuity set for the law of \(X_\infty\), so
\cref{thm:p15-weak-portmanteau}(iv) does not apply to them, and indeed
\(\mathbb P(X_n\in G)\not\to\mathbb P(X_\infty\in G)\).
\item By contrast, if \(x_n\to x_\infty\) with \(x_\infty\in G\), then
\(x_n\in G\) for all sufficiently large \(n\) and
\(\mathbb P(X_n\in G)\to1=\mathbb P(X_\infty\in G)\).
\end{enumerate}
\end{hbtheorem}

The remark is a mnemonic for the directions of the inequalities: mass
sitting inside an open set can leave it in the limit, through the
boundary, but cannot enter it; for a closed set the reverse holds. The
instance \(G=(0,\infty)\), \(x_n=1/n\), \(x_\infty=0\) is the first
item of \cref{thm:p15-weak-convergence-boundaries}.

\begin{hbexample}{def:p15-w02c-p15-004-two-sided}{A boundary point approached from both sides}
Let \(G=\mathbb R\setminus\{0\}\) and \(x_n=(-1)^n/n\). The set \(G\)
is open, because the complement \(\{0\}\) is closed. Every interval
around \(0\) contains nonzero points, so \(\overline G=\mathbb R\);
since \(G^\circ=G\), \(\partial G=\{0\}\). The points \(x_n\) lie in
\(G\), alternate in sign, and converge to \(x_\infty=0\). By
\cref{thm:p15-w02c-p15-004-boundary-approach},
\(\delta_{x_n}\Rightarrow\delta_0\), every \(\delta_{x_n}\) gives
\(G\) mass one while \(\delta_0\) gives it mass zero, and the closed set
\(K=\{0\}\) receives mass zero along the sequence and mass one in the
limit. The distribution functions show the same thing directly: for
\(x\neq0\), \(F_n(x)=\mathbf1\{x_n\leq x\}\) equals \(\mathbf1\{x>0\}\)
once \(1/n<|x|\), which is the limit distribution function at \(x\);
at the single discontinuity point \(x=0\) the values
\(F_n(0)=\mathbf1\{n\hbox{ odd}\}\) do not converge, which
\cref{def:p15-weak-law-convergence} permits.
\end{hbexample}

\section{Nice spaces, Borel real codings and the extension theorem}
\label{sec:p15-w02c-p15-004-nice}

\Cref{thm:p15-kolmogorov-countable} extends consistent laws on
\(\mathbb R^n\) to the space of real sequences, and
\cref{thm:p15-coded-space-extension} transports that extension to
factors carrying a Borel real coding in the sense of
\cref{def:p15-measurable-real-coding}. Durrett reaches the same point
through a class of measurable spaces he calls nice
\cite[p.~53]{DurrettPTEAuthor2019}. His definition does not ask that
the coding have a Borel image, while the handbook's coding does. This
section states his definition and his Theorem~2.1.22 in his own terms,
proves the exact relation between nice maps and Borel real codings,
reports his statement that the extension theorem fails for arbitrary
measurable spaces together with the counterexamples he attributes, and
sets his nice-space generalization against the handbook's Borel-image
version. The last result shows that the Borel-image clause is not a
convenience: without it the extension can fail.

\begin{hbdefinition}{def:p15-w02c-p15-004-nice}{Durrett's nice measurable spaces}
A measurable space \((S,\mathcal S)\) is \emph{nice} when there is a
one-to-one map \(\varphi\) from \(S\) into \(\mathbb R\) such that
\(\varphi\) and \(\varphi^{-1}\) are both measurable
\cite[p.~53]{DurrettPTEAuthor2019}. Typed in full: \(\varphi:S\to\mathbb R\)
is injective; \(\varphi^{-1}(B)\in\mathcal S\) for every
\(B\in\mathcal B(\mathbb R)\); and the inverse map
\(\varphi^{-1}:\varphi(S)\to S\), defined on the image, is measurable from
\(\varphi(S)\) with its trace sigma-algebra
\(\{\varphi(S)\cap B:B\in\mathcal B(\mathbb R)\}\) to \((S,\mathcal S)\),
that is, \(\varphi(A)\) belongs to that trace sigma-algebra for every
\(A\in\mathcal S\). Such a \(\varphi\) is called a \emph{nice map} here.
No clause requires \(\varphi(S)\) to be a Borel set.

Durrett adds that these spaces are often called \emph{standard Borel
spaces}, a name he declines because, as he puts it, ``we already have too
many things named after Borel'' \cite[p.~53]{DurrettPTEAuthor2019}. He
introduces the class as one adequate for all his later results and on
which the generalization of Kolmogorov's theorem is trivial, and says that
his next result shows that most spaces arising in applications are nice
\cite[p.~53]{DurrettPTEAuthor2019}. In this handbook the phrase
\emph{standard Borel space} keeps the meaning fixed in
\cref{def:p15-measurable-real-coding}: a space with a Borel real coding.
\Cref{thm:p15-w02c-p15-004-nice-coding} gives the exact relation of the
two notions, and \cref{thm:p15-w02c-p15-004-nice-extension} shows that
they differ.
\index{nice measurable space}\index{standard Borel space!nice space}
\end{hbdefinition}

\begin{hbtheorem}{thm:p15-w02c-p15-004-nice-coding}{proved=annex:p15-w02c-p15-004-nice-coding}{Nice maps and Borel real codings}
Let \((S,\mathcal S)\) be a measurable space and \(\varphi:S\to\mathbb R\).
\begin{enumerate}
\item If \(\varphi\) is a Borel real coding of \((S,\mathcal S)\) onto a
      Borel set \(E\), then \(\varphi\) is a nice map and \(E=\varphi(S)\).
\item If \(\varphi\) is a nice map, then \(\varphi\) is a Borel real coding
      of \((S,\mathcal S)\) onto \(\varphi(S)\) if and only if
      \(\varphi(S)\in\mathcal B(\mathbb R)\).
\item If \(\varphi\) is a nice map, then \(\varphi(A)\in\mathcal B(\mathbb R)\)
      for every \(A\in\mathcal S\) if and only if
      \(\varphi(S)\in\mathcal B(\mathbb R)\).
\item For every set \(A\subseteq\mathbb R\), equipped with its trace
      sigma-algebra \(\{A\cap B:B\in\mathcal B(\mathbb R)\}\), the
      inclusion map \(A\to\mathbb R\) is a nice map; it is a Borel real
      coding exactly when \(A\in\mathcal B(\mathbb R)\).
\end{enumerate}
\end{hbtheorem}

So a space with a Borel real coding is nice, and a nice space is a
standard Borel space in the handbook's sense as soon as some nice map
for it has a Borel image. Part~(iii) shows what the Borel-image clause
buys: it is exactly the condition under which a nice map carries every
measurable set to a Borel subset of the line.

\begin{hbtheorem}{thm:p15-w02c-p15-004-durrett-2122}{proved=annex:p15-w02c-p15-004-durrett-2122}{Durrett's Theorem~2.1.22 in the source's terms}
If \(S\) is a Borel subset of a complete separable metric space \(M\),
and \(\mathcal S\) is the collection of Borel subsets of \(S\), then
\((S,\mathcal S)\) is nice \cite[Theorem~2.1.22, p.~54]{DurrettPTEAuthor2019}.
More exactly:
\begin{enumerate}
\item The Borel sigma-algebra of \(S\) for the restricted metric, the
      trace sigma-algebra \(\{S\cap B:B\in\mathcal B(M)\}\) of
      \cref{def:p15-measurable-real-coding}, and the family
      \(\{B\in\mathcal B(M):B\subseteq S\}\) coincide; so the source's
      \(\mathcal S\) is the same sigma-algebra on either reading of
      ``Borel subsets of \(S\)''.
\item The coding \(\phi\) of \cref{thm:p15-metric-borel-representation}
      is a nice map for \((S,\mathcal S)\) whose image is a Borel subset
      of \(\mathbb R\).
\end{enumerate}
\end{hbtheorem}

Durrett's proof composes the binary interleaving of \([0,1)^{\mathbb N_+}\)
with the distance map \(x\mapsto(d(x,q_1),d(x,q_2),\ldots)\) and says that
the composite gives the required mapping; he then warns that the proof is
light on details \cite[p.~54]{DurrettPTEAuthor2019}. The measurability of
the inverse and the Borel image are proved in
\cref{thm:p15-sequence-borel-coding,thm:p15-metric-borel-representation};
part~(ii) above reads that result in the vocabulary of nice maps.
The source's applications claim is thereby established in the following
exact form: every Borel subset of a complete separable metric space, with
its Borel sets, is nice, and is nice through a map with Borel image. The
handbook proves the hypotheses directly for \([0,1)^{\mathbb N_+}\) with the
metric of \cref{thm:p15-sequence-borel-coding}; no further catalogue of
applied spaces is asserted here.

\begin{hbwarning}{def:p15-w02c-p15-004-extension-failure}{Durrett's failure claim and attributions}
Durrett states that his Theorem~2.1.21, Kolmogorov's extension theorem for
consistent probability measures on \((\mathbb R^n,\mathcal R^n)\), is not
valid for arbitrary measurable spaces \((S,\mathcal S)\)
\cite[p.~53]{DurrettPTEAuthor2019}. He attributes the first example, on an
infinite product \(\Omega_1\times\Omega_2\times\cdots\) of different
spaces, to Andersen and Jessen (1948), and for it refers the reader to
Halmos (1950), page~214, or to Neveu (1965), page~84; for an example in
which all the spaces \(\Omega_i\) are the same, he refers to Wegner (1973)
\cite[p.~53]{DurrettPTEAuthor2019}. Those four works are reported here
only as the source names them, with the source's page references; their
constructions are not reproduced, and the example in
\cref{thm:p15-w02c-p15-004-nice-extension}(iii) is not claimed to be any
of them. It shares the shape of the first one: its factor spaces are
different.
\index{Kolmogorov extension theorem!failure for arbitrary spaces}
\end{hbwarning}

The generalization that Durrett calls trivial is, read with his definition
of nice, the following statement: \emph{if each \((S_i,\mathcal S_i)\),
\(i\geq1\), is nice and probability laws \(\nu_n\) on the finite products
\(\prod_{i\leq n}S_i\) are consistent, then there is a unique probability
on \(\prod_{i\geq1}S_i\) with those finite-dimensional laws.} The source
gives no formal statement or proof of it on page~53. The next theorem
separates what holds from what does not.

\begin{hbtheorem}{thm:p15-w02c-p15-004-nice-extension}{proved=annex:p15-w02c-p15-004-nice-extension}{Extension over nice factors: the Borel-image clause is needed}
For measurable spaces \((S_i,\mathcal S_i)\), \(i\geq1\), give each finite
product \(\prod_{i\leq n}S_i\) and the infinite product
\(\prod_{i\geq1}S_i\) the sigma-algebra generated by the coordinate sets
\(\{x:x_i\in A\}\), \(A\in\mathcal S_i\). Call probability laws
\(\nu_n\) on \(\prod_{i\leq n}S_i\) \emph{consistent} when
\(\nu_{n+1}(E\times S_{n+1})=\nu_n(E)\) for every \(n\) and every
measurable \(E\subseteq\prod_{i\leq n}S_i\). A probability \(P\) on
\(\prod_{i\geq1}S_i\) \emph{has the finite-dimensional laws} \(\nu_n\)
when \(P(\{x:(x_1,\ldots,x_n)\in E\})=\nu_n(E)\) for every \(n\) and
every measurable \(E\subseteq\prod_{i\leq n}S_i\).
\begin{enumerate}
\item \emph{Uniqueness holds for arbitrary factors.} Two probabilities on
      \(\prod_{i\geq1}S_i\) with the same finite-dimensional laws are
      equal.
\item \emph{Existence with Borel images.} If each \((S_i,\mathcal S_i)\)
      has a nice map with Borel image, then every consistent sequence
      \((\nu_n)\) has exactly one probability on \(\prod_{i\geq1}S_i\)
      with the finite-dimensional laws \(\nu_n\).
\item \emph{Failure for nice factors without the Borel-image clause.}
      Assume the Axiom of Choice of \cref{sec:set-prerequisites-06}.
      There are sets \(A_1\supseteq A_2\supseteq\cdots\) contained in
      \([0,1)\) with \(\bigcap_nA_n=\varnothing\) such that, with
      \(S_i=A_i\) and \(\mathcal S_i\) its trace sigma-algebra, each
      inclusion \(A_i\to\mathbb R\) is a nice map, there is a consistent
      sequence \((\nu_n)\), and no probability on
      \(\prod_{i\geq1}S_i\) has the finite-dimensional laws \(\nu_n\).
      In particular at least one of these nice spaces has no Borel real
      coding, so the class of nice spaces is strictly larger than the
      class of standard Borel spaces in the handbook's sense.
\end{enumerate}
\end{hbtheorem}

Part~(iii) is where the handbook's narrowing stands against the source.
Durrett's claim that the generalization is trivial holds, with proof, when
each nice map has a Borel image, which is the form proved in
\cref{thm:p15-coded-space-extension}; under the literal definition of
\cref{def:p15-w02c-p15-004-nice} the existence half of the statement is
false, while the uniqueness half survives by part~(i). For every space the
source proves nice, \cref{thm:p15-w02c-p15-004-durrett-2122}(ii) supplies a
nice map with Borel image, so the extension is available on those spaces
by part~(ii).

\begin{hbexample}{def:p15-w02c-p15-004-nice-cases}{Three nice maps}
The identity of \(\mathbb R\), with \(\mathcal S=\mathcal B(\mathbb R)\), is
a nice map with Borel image, hence a Borel real coding by
\cref{thm:p15-w02c-p15-004-nice-coding}(ii). The source's map
\(\phi_0\) of \cref{def:p15-binary-sequence-code} is a nice map for
\([0,1)^{\mathbb N_+}\) with its product sigma-algebra, and its image is
the Borel set \(E_0\) by \cref{thm:p15-sequence-borel-coding}. Each set
\(A_n\) of \cref{thm:p15-w02c-p15-004-nice-extension}(iii) is nice through
its inclusion, but that inclusion is not a Borel real coding, because
\(A_n\) is not a Borel set. Indeed, the proof shows that every Borel set
containing \(A_n\) has Lebesgue measure one in \([0,1)\), and that
countably many rotations of the complement \([0,1)\setminus A_n\) cover
\([0,1)\); if \(A_n\) were Borel, its complement would be a Borel set of
measure zero whose countably many rotations, each of measure zero, cover
a set of measure one.
\end{hbexample}

\section{Cardinality of Borel sets in complete separable metric spaces}
\label{sec:p15-w02c-p15-004-borel-cardinality}

After the proof of his Theorem~2.1.22, Durrett notes that the analysis
in Section~13.1 of Dudley's \emph{Real Analysis and Probability} (1989)
has a consequence for cardinality: a Borel subset of a complete
separable metric space is finite, countably infinite, or of the
cardinality of the real numbers \cite[p.~54]{DurrettPTEAuthor2019}. The
source gives no proof, and \cref{def:p15-countable-product-next} records
that the injective coding of \cref{thm:p15-metric-borel-representation}
does not by itself supply one. This section proves the statement. The
coding reduces it to Borel subsets of \(\mathbb R\); Borel subsets of
\(\mathbb R\) are continuous images of closed subsets of the space of
natural-number sequences; and an uncountable such image receives an
injection of the binary sequences. The last step is a form of the
perfect set theorem for Suslin sets; Moschovakis proves it by a
transfinite derivative of trees
\cite[2C.2, pp.~59--61; 2C.3, p.~61]{Moschovakis2009DSTAuthor}, while
the argument here removes the countably many basic pieces with countable
image in a single step. Dudley's own treatment is not used.

\begin{hbdefinition}{def:p15-w02c-p15-004-baire-space}{Sequence space, closed sets, Suslin images}
Let \(\mathbb N=\{0,1,2,\ldots\}\) and \(\mathcal N=\mathbb N^{\mathbb N}\).
For a finite sequence \(s=(s_0,\ldots,s_{k-1})\) of natural numbers,
of length \(|s|=k\), put
\(\mathcal N_s=\{\alpha\in\mathcal N:\alpha(i)=s_i\hbox{ for }i<k\}\),
and write \(\alpha|k=(\alpha(0),\ldots,\alpha(k-1))\). A set
\(F\subseteq\mathcal N\) is \emph{closed} when for every
\(\alpha\notin F\) there is \(k\) with \(\mathcal N_{\alpha|k}\cap F=\varnothing\).
A function \(f:F\to\mathbb R\) is \emph{continuous} when for every
\(\alpha\in F\) and \(\varepsilon>0\) there is \(k\) with
\(|f(\beta)-f(\alpha)|<\varepsilon\) for all
\(\beta\in F\cap\mathcal N_{\alpha|k}\). A set \(A\subseteq\mathbb R\)
is a \emph{Suslin image} when \(A=\varnothing\) or \(A=f[F]\) for some
closed \(F\subseteq\mathcal N\) and continuous \(f:F\to\mathbb R\).
Write \(2^{\mathbb N}=\{0,1\}^{\mathbb N}\) for the binary sequences.
A set is \emph{countable} when it is finite or countably infinite.
\index{Baire space}\index{Suslin image}
\end{hbdefinition}

\begin{hblemma}{thm:p15-w02c-p15-004-reals-continuum}{proved=annex:p15-w02c-p15-004-reals}{The real line and the binary sequences are equinumerous}
There is a bijection \(\mathbb R\to2^{\mathbb N}\). In particular
\(\mathbb R\) is not countable.
\end{hblemma}

\begin{hblemma}{thm:p15-w02c-p15-004-borel-suslin}{proved=annex:p15-w02c-p15-004-borel-suslin}{Borel subsets of the line are Suslin images}
Every closed subset of \(\mathbb R\) is a Suslin image. Suslin images
are closed under countable unions and countable intersections.
Consequently every Borel subset of \(\mathbb R\) is a Suslin image.
\end{hblemma}

\begin{hblemma}{thm:p15-w02c-p15-004-condensation}{proved=annex:p15-w02c-p15-004-cond}{An uncountable Suslin image contains a copy of the binary sequences}
If \(F\subseteq\mathcal N\) is closed, \(f:F\to\mathbb R\) is
continuous and \(f[F]\) is uncountable, then there is an injection
\(2^{\mathbb N}\to f[F]\).
\end{hblemma}

\begin{hbtheorem}{thm:p15-w02c-p15-004-borel-cardinality}{proved=annex:p15-w02c-p15-004-borel-cardinality}{Durrett, p.~54: Borel sets satisfy the continuum dichotomy}
Let \((M,\rho)\) be a complete separable metric space and let
\(S\subseteq M\) be Borel. Then exactly one of the following holds:
\(S\) is finite; \(S\) is countably infinite; there is a bijection
\(S\to\mathbb R\).
\end{hbtheorem}

The hypothesis that \(S\) is Borel cannot be dropped from the statement
as proved. Whether every uncountable subset of \(\mathbb R\) is in
bijection with \(\mathbb R\) is the continuum hypothesis, and
\cref{thm:p06-many-cohen-continuum} produces models in which
\(2^{\aleph_0}>\aleph_1\). In such a model a subset of \(\mathbb R\) of
cardinality \(\aleph_1\) is uncountable and not in bijection with
\(\mathbb R\), so by \cref{thm:p15-w02c-p15-004-borel-cardinality} it is
not Borel. The theorem is therefore a statement about the Borel sets,
not a proof of the continuum hypothesis.

\begin{hbexample}{def:p15-w02c-p15-004-irrationals}{The three cases on the line}
The Borel set \(\{0,1\}\) is finite. The set \(\mathbb Q\) is Borel, as a
countable union of closed singletons, and is countably infinite. The
irrationals \(\mathbb R\setminus\mathbb Q\) form a Borel set, being the
complement of a Borel set. They contain no interval, since every
nondegenerate interval contains a rational. They are uncountable: if they
were countable, \(\mathbb R\) would be a union of two countable sets and
so countable, contrary to \cref{thm:p15-w02c-p15-004-reals-continuum}.
\Cref{thm:p15-w02c-p15-004-borel-cardinality} therefore yields a bijection
of \(\mathbb R\setminus\mathbb Q\) with \(\mathbb R\), although no interval
lies inside it. The proof supplies the injection of \(2^{\mathbb N}\) into
this set explicitly through
\cref{thm:p15-w02c-p15-004-borel-suslin,thm:p15-w02c-p15-004-condensation}.
\end{hbexample}

\section{The birthday approximation and the source's density notation}
\label{sec:p15-w02c-p15-004-examples}

\Cref{thm:p15-weak-first-examples} proves the limit law of the birthday
time, and \cref{thm:p15-weak-central-order-statistic} proves the density
and the normal limit of the scaled median. Durrett's Example~3.2.5 also
says what the birthday probability means and how far its limit
approximation is from the exact value, and his Example~3.2.6 writes
densities in a notation of its own
\cite[pp.~117--118]{DurrettPTEAuthor2019}. This section states both
points in the source's terms and proves the numerical comparison.

\begin{hbexample}{def:p15-w02c-p15-004-birthday-reading}{The reading of Durrett Example~3.2.5}
The source's model takes independent variables \(X_1,X_2,\ldots\),
each uniformly distributed on \(\{1,\ldots,N\}\), the time of the first repeat is
\[
 T_N=\min\{n:X_n=X_m\text{ for some }m<n\},
\]
and
\[
 P(T_N>n)=\prod_{m=2}^{n}\left(1-\frac{m-1}{N}\right)
\]
\cite[Example~3.2.5, p.~117]{DurrettPTEAuthor2019}. The substitution
\(j=m-1\) turns this into the product
\(\prod_{j=1}^{n-1}(1-j/N)\) of \cref{thm:p15-weak-first-examples}, for
\(1\leq n\leq N+1\). The event \(\{T_N>n\}\) is the event that
\(X_1,\ldots,X_n\) are pairwise distinct.

\emph{Reading.} Take \(N=365\) and let \(X_1,\ldots,X_n\) be the
birthdays of the \(n\) members of a group, under the source's
assumption that all birthdays are equally likely, together with the
independence built into the model. Then \(P(T_{365}>n)\) is the
probability that no two members of the group share a birthday. The
source states this as the probability that two people in a group of
size \(n\) do not have the same birthday \cite[p.~117]{DurrettPTEAuthor2019};
the event meant is that no pair of members shares one.

\emph{The source's approximation.} From the limit
\[
 P(T_N/N^{1/2}>x)\to\exp(-x^2/2)\qquad(x\geq0),
\]
the source takes \(N=365\), computes \(22/\sqrt{365}=1.1515\) and
\((1.1515)^2/2=0.6630\), and concludes \(P(T_{365}>22)\approx
e^{-0.6630}\approx0.515\). It then compares: the approximation is
\(2\%\) smaller than the true probability \(0.524\)
\cite[p.~117]{DurrettPTEAuthor2019}. \Cref{thm:p15-w02c-p15-004-birthday-error}
proves the comparison: the shortfall is between \(1.71\) and \(1.74\)
percent of the exact value, which rounds to the source's \(2\%\).

\emph{Limitation.} The limit theorem holds for fixed \(x\) as
\(N\to\infty\) and carries no error bound at a fixed \(N\); the
comparison below is an evaluation at the single pair \(N=365\),
\(n=22\), not a rate of convergence. The numbers describe the uniform
independent model only.
\end{hbexample}

\begin{hbtheorem}{thm:p15-w02c-p15-004-birthday-error}{proved=annex:p15-w02c-p15-004-birthday-error}{Error of the birthday approximation}
Let
\[
 p=P(T_{365}>22)=\prod_{j=1}^{21}(1-j/365),\qquad
 a=e^{-22^2/(2\cdot365)}=e^{-484/730}.
\]
Then \(a<p\) and
\[
 0.0171<\frac{p-a}{p}<0.0174 .
\]
The same two bounds hold with \(a\) replaced by the source's value
\(e^{-0.6630}\), computed from the rounded exponent.
\end{hbtheorem}

\Cref{thm:p15-weak-first-examples} records the decimal values
\(p\approx0.524305\) and \(a\approx0.515296\); the theorem above gives
the relative comparison without relying on those decimals.

\begin{hbconvention}{def:p15-w02c-p15-004-density-notation}{Durrett's notation \(P(Y_n=y)\) for a density}
After computing the density of \(Y_n=2(V_{n+1}-1/2)\sqrt{2n}\) in
Example~3.2.6, Durrett fixes a convention:
``Here and in what follows we write \(P(Y_n=y)\) for the density
function of \(Y_n\)'' \cite[p.~118]{DurrettPTEAuthor2019}. In the
source, \(P(Y_n=y)\) therefore stands for the value at \(y\) of the
density written \(f_{Y_n}(y)\) in
\cref{thm:p15-weak-central-order-statistic}, and the source's limit
statement for the density is the pointwise limit
\(f_{Y_n}(y)\to(2\pi)^{-1/2}\exp(-y^2/2)\) proved there.

This handbook does not adopt the convention. Densities are written
\(f_{Y_n}\), and \(P(Y_n=y)\), whenever it occurs here, is the
probability of the event \(\{Y_n=y\}\). Read in that sense the symbol is
zero for every \(y\), while the density is not
(\cref{thm:p15-w02c-p15-004-density-reading}), so the two readings
cannot be interchanged. In the same display of the source, the first
factor of the density, \(\binom{2n}{n}2^{-2n}\), is identified as
\(P(S_{2n}=0)\) for a simple random walk; that symbol is a probability
of an event in the ordinary sense, and equals \(C_n\) of
\cref{thm:p15-weak-elementary-asymptotics}, with \(S_{2n}\) as in
\cref{thm:p15-weak-rademacher-example}.
\end{hbconvention}

\begin{hblemma}{thm:p15-w02c-p15-004-density-reading}{proved=annex:p15-w02c-p15-004-dens}{The two readings of \(P(Y_n=y)\) differ}
Let \(n\geq1\) and let \(Y_n\) be as in
\cref{thm:p15-weak-central-order-statistic}.
\begin{enumerate}
\item For every real \(y\), the event \(\{Y_n=y\}\) has probability
zero.
\item \(f_{Y_n}(0)=C_n\,\frac{2n+1}{2n}\sqrt{n/2}>0\).
\item If \(\varepsilon_1,\ldots,\varepsilon_{2n}\) are independent with
probabilities \(1/2\) at \(1\) and \(-1\), and
\(S_{2n}=\sum_{i=1}^{2n}\varepsilon_i\), then
\(\mathbb P(S_{2n}=0)=\binom{2n}{n}2^{-2n}=C_n\).
\end{enumerate}
\end{hblemma}

\section{Relationships and sources}
\label{sec:p15-w02c-p15-004-relations}

\Cref{thm:p15-w02c-p15-004-topological-weak}(ii) is
\cref{thm:p15-weak-bounded-continuous} read in the general vocabulary;
part~(iii) is the continuous case of the continuous mapping theorem
\cref{thm:p15-weak-continuous-mapping}, now on general spaces and without
an exceptional set. The proof of part~(iv) uses only the uniqueness
theorem for measures that agree on a generating pi-system,
\cref{thm:p15-measure-uniqueness}. The point-mass results are applications
of the bounded-continuous characterization and of the portmanteau
theorem \cref{thm:p15-weak-portmanteau}; they extend the single instance
in \cref{thm:p15-weak-convergence-boundaries} to every open set and every
boundary approach.

\Cref{thm:p15-w02c-p15-004-nice-coding,thm:p15-w02c-p15-004-durrett-2122}
connect the source's nice spaces with \cref{def:p15-measurable-real-coding}
and \cref{thm:p15-metric-borel-representation};
\cref{thm:p15-w02c-p15-004-nice-extension}(ii) is
\cref{thm:p15-coded-space-extension} restated for nice maps with Borel
images, and part~(iii) shows that the restriction to Borel images in that
theorem cannot be removed. The uniqueness in part~(i) is the product-space
form of \cref{thm:p15-measure-uniqueness}, and the measure used in the
counterexample is Lebesgue measure of \cref{thm:p15-lebesgue-unit-interval}.

\Cref{thm:p15-w02c-p15-004-borel-cardinality} completes the remark
recorded at \cref{def:p15-countable-product-next}: the injective coding
of \cref{thm:p15-metric-borel-representation} transfers the question to
the line, and \cref{thm:p15-w02c-p15-004-borel-suslin,thm:p15-w02c-p15-004-condensation}
answer it there. The perfect set property of analytic sets in Polish
spaces is surveyed, with a cited proof, in
\cref{thm:p06-analytic-regularity}; the two lemmas here prove directly
the case needed, for Borel subsets of the line.

\Cref{def:p15-w02c-p15-004-birthday-reading,thm:p15-w02c-p15-004-birthday-error}
extend the birthday item of \cref{thm:p15-weak-first-examples}, whose
product-limit lemma \cref{thm:p15-weak-elementary-asymptotics} gives the
limit used by the approximation;
\cref{def:p15-w02c-p15-004-density-notation,thm:p15-w02c-p15-004-density-reading}
fix how the source's notation in Example~3.2.6 maps onto
\cref{thm:p15-weak-central-order-statistic}, whose total-variation
conclusion comes from Scheff\'e's theorem
\cref{thm:p15-scheffe-density-limit}.

Sources: the test-function remark and Theorem~3.2.9 are on page~119
\cite[p.~119]{DurrettPTEAuthor2019}; Theorem~3.2.11 and the boundary remark are
on page~120 \cite[p.~120]{DurrettPTEAuthor2019}; Theorem~2.1.21, the
failure statement with its attributions, the definition of nice spaces
and the remarks on their name and on applications are on page~53
\cite[p.~53]{DurrettPTEAuthor2019}, and Theorem~2.1.22 with its proof is
on page~54 \cite[Theorem~2.1.22, p.~54]{DurrettPTEAuthor2019}; the cardinality
remark, with its pointer to Dudley's Section~13.1, is on page~54
\cite[p.~54]{DurrettPTEAuthor2019}. Example~3.2.5, with its reading, its arithmetic and its
comparison with the true probability, is on page~117
\cite[Example~3.2.5, p.~117]{DurrettPTEAuthor2019}, and the density
convention of Example~3.2.6 is on page~118
\cite[p.~118]{DurrettPTEAuthor2019}. For the tree-derivative proof of the
perfect set theorem for Suslin sets see
\cite[2C.2, pp.~59--61]{Moschovakis2009DSTAuthor}.
